% Introduction. Numbers: ../FINDINGS.md.
\section{Introduction}\label{sec:intro}

Retractions have become far more frequent. In the Retraction Watch database, 13,210 papers have a first retraction
notice dated 2023, more than in any earlier year, and papers retracted for paper-mill reasons rose from 246 in 2020 to
6,782 in 2023
\citep{rittman2025rw,candal2022papermills}. Because retractions are slow, often years after publication, editors and integrity
teams look for early signals. One signal is already in use: the ``Feet of Clay'' detector of the Problematic Paper
Screener flags articles whose reference lists include retracted papers \citep{cabanac2022pps}. Network studies support
the idea. Retracted papers cluster in citation networks \citep{feng2024molbio,candal2024references}, and in a
retrospective analysis of more than 50,000 Web of Science retractions, articles citing retracted works had
substantially higher subsequent retraction rates \citep{lou2026beyond}.

That association has not been tested in the way a screening tool is used: with information frozen at one moment and
outcomes observed afterwards. A retrospective association can also arise from journals alone. Retractions arrive in
waves within journals and special issues, and papers in the same venue cite one another, so a paper that cites a
retracted paper may simply sit in a journal that is being cleaned up. Machine-learning studies of retraction
prediction so far have used text and metadata, not citation links, and random rather than temporal splits
\citep{fletcher2025predicting}.

This article tests the signal prospectively on open data. All predictors are computed on the January 2024
OpenCitations graph \citep{arjomandi2024dataset} with the retractions known at that time, frozen with a hash, and then
compared with retractions issued from February 2024 to October 2026, using the Retraction Watch database
\citep{rittman2025rw}. The scores, metrics and decision rule were pre-registered. The contributions are:
\begin{itemize}
\item a prospective, pre-registered test of whether citing (or being cited by) a retracted paper predicts later
retraction, on 10.5 million biomedical papers;
\item a separation of the neighbourhood signal from journal-level clustering, using strata of journal and publication
year;
\item an assessment of its value as a screen (precision and lift) against a journal baseline and a model fitted on an
earlier window;
\item supporting measurements of community concentration, NIH-funded exposure to retracted work and citations around
retraction, all from open data, with code and run records.
\end{itemize}
